\documentclass[10pt,a4paper]{amsart}
\usepackage{amsmath,amssymb,amsthm,graphicx}
\usepackage{verbatim}
\newenvironment{reference}{\comment}{\endcomment}
%\newenvironment{reference}{}{}
\newenvironment{slogan}{\comment}{\endcomment}
\newenvironment{history}{\comment}{\endcomment}

% For commutative diagrams we use Xy-pic
\usepackage[all]{xy}

% We use 2cell for 2-commutative diagrams.
\xyoption{2cell}
% Theorem environments.
%
\theoremstyle{plain}
\newtheorem{theorem}[subsection]{Theorem}
\newtheorem{proposition}[subsection]{Proposition}
\newtheorem{lemma}[subsection]{Lemma}

\theoremstyle{definition}
\newtheorem{definition}[subsection]{Definition}
\newtheorem{example}[subsection]{Example}
\newtheorem{exercise}[subsection]{Exercise}
\newtheorem{situation}[subsection]{Situation}

\theoremstyle{remark}
\newtheorem{remark}[subsection]{Remark}
\newtheorem{remarks}[subsection]{Remarks}

\numberwithin{equation}{subsection}

% Macros
%
\def\lim{\mathop{\mathrm{lim}}\nolimits}
\def\colim{\mathop{\mathrm{colim}}\nolimits}
\def\Spec{\mathop{\mathrm{Spec}}}
\def\Hom{\mathop{\mathrm{Hom}}\nolimits}
\def\Ext{\mathop{\mathrm{Ext}}\nolimits}
\def\SheafHom{\mathop{\mathcal{H}\!\mathit{om}}\nolimits}
\def\SheafExt{\mathop{\mathcal{E}\!\mathit{xt}}\nolimits}
\def\Sch{\mathit{Sch}}
\def\Mor{\mathop{\mathrm{Mor}}\nolimits}
\def\Ob{\mathop{\mathrm{Ob}}\nolimits}
\def\Sh{\mathop{\mathit{Sh}}\nolimits}
\def\NL{\mathop{N\!L}\nolimits}
\def\CH{\mathop{\mathrm{CH}}\nolimits}
\def\proetale{{pro\text{-}\acute{e}tale}}
\def\etale{{\acute{e}tale}}
\def\QCoh{\mathit{QCoh}}
\def\Ker{\mathop{\mathrm{Ker}}}
\def\Im{\mathop{\mathrm{Im}}}
\def\Coker{\mathop{\mathrm{Coker}}}
\def\Coim{\mathop{\mathrm{Coim}}}

% Boxtimes
%
\DeclareMathSymbol{\boxtimes}{\mathbin}{AMSa}{"02}

%
% Macros for moduli stacks/spaces
%
\def\QCohstack{\mathcal{QC}\!\mathit{oh}}
\def\Cohstack{\mathcal{C}\!\mathit{oh}}
\def\Spacesstack{\mathcal{S}\!\mathit{paces}}
\def\Quotfunctor{\mathrm{Quot}}
\def\Hilbfunctor{\mathrm{Hilb}}
\def\Curvesstack{\mathcal{C}\!\mathit{urves}}
\def\Polarizedstack{\mathcal{P}\!\mathit{olarized}}
\def\Complexesstack{\mathcal{C}\!\mathit{omplexes}}
% \Pic is the operator that assigns to X its picard group, usage \Pic(X)
% \Picardstack_{X/B} denotes the Picard stack of X over B
% \Picardfunctor_{X/B} denotes the Picard functor of X over B
\def\Pic{\mathop{\mathrm{Pic}}\nolimits}
\def\Picardstack{\mathcal{P}\!\mathit{ic}}
\def\Picardfunctor{\mathrm{Pic}}
\def\Deformationcategory{\mathcal{D}\!\mathit{ef}}
\usepackage{cleveref}
\crefname{originalequation}{original equation}{original equations}
\Crefname{originalequation}{Original equation}{Original equations}
\makeatletter
\newlabel{equation-epsilon}{{08H7}{source}{}{sourcepointer-1}{}}
\newlabel{equation-epsilon@cref}{{[originalequation][0][]08H7}{[1][0][]source}{}{}{}}
\makeatother
\crefname{originallemma}{original lemma}{original lemmas}
\Crefname{originallemma}{Original lemma}{Original lemmas}
\makeatletter
\newlabel{morphisms-lemma-etale-open}{{03WT}{source}{}{sourcepointer-2}{}}
\newlabel{morphisms-lemma-etale-open@cref}{{[originallemma][0][]03WT}{[1][0][]source}{}{}{}}
\makeatother
\crefname{originallemma}{original lemma}{original lemmas}
\Crefname{originallemma}{Original lemma}{Original lemmas}
\makeatletter
\newlabel{sites-modules-lemma-local-pullback}{{03DO}{source}{}{sourcepointer-3}{}}
\newlabel{sites-modules-lemma-local-pullback@cref}{{[originallemma][0][]03DO}{[1][0][]source}{}{}{}}
\makeatother
\title{Finite type of a module descends from the small etale site}
\author{The Stacks Project Authors}
\date{stacks-project:spaces-perfect.tex:local-snapshot-20261012; GNU FDL 1.2 or later}
\begin{document}
\maketitle
\noindent\textit{Editorial note.} All delivered mathematical text is the author's own, copied from the pinned Stacks Project chapter file spaces-perfect.tex; no statement or proof step was rewritten, reordered or removed.\par\medskip
\noindent\textit{Editorial note.} The counted proof is the author's complete argument: the proof is complete in both directions\par\medskip
\noindent\textit{Editorial note.} The label lemma-descend-finite-type names the counted statement, and the proof environment immediately following it in the source is delivered verbatim as its proof.\par\medskip
\noindent\textit{Editorial note.} The counted statement is carried by the author's own lemma environment in the chapter source.\par\medskip
\noindent\textit{Allowed standard prerequisites (editorial).}\begin{itemize}
\item the localization functor epsilon from the small etale site to the Zariski site and the description of epsilon\textasciicircum{}*F as epsilon\textasciicircum{}\{-1\}F tensor O\_etale
\item generation by finitely many sections on an etale covering
\item construction of the inverse image functor and descent of sections after shrinking
\item faithful flatness of etale morphisms of schemes on stalks
\end{itemize}\medskip
\ifdefined\hypertarget\hypertarget{sourcepointer-1}{}\fi
\noindent\textit{Source reference.} Original equation 08H7: sources/stacks-spaces-perfect/pinned-original-source.tex, source lines 320--320. Prerequisite pointer invoked at proof line(s) 3089; it names a standard fact of the same chapter that the counted argument uses, and every step of the counted proof built on it is retained verbatim in the delivered body.\par\medskip
\ifdefined\hypertarget\hypertarget{sourcepointer-2}{}\fi
\noindent\textit{Source reference.} Original lemma 03WT: sources/stacks-chow/tags.txt, source lines 5124--5124. Prerequisite pointer invoked at proof line(s) 3113; it names a standard fact of a sibling chapter that the counted argument uses, and every step of the counted proof built on it is retained verbatim in the delivered body.\par\medskip
\ifdefined\hypertarget\hypertarget{sourcepointer-3}{}\fi
\noindent\textit{Source reference.} Original lemma 03DO: sources/stacks-chow/tags.txt, source lines 4427--4427. Prerequisite pointer invoked at proof line(s) 3094; it names a standard fact of a sibling chapter that the counted argument uses, and every step of the counted proof built on it is retained verbatim in the delivered body.\par\medskip
\begin{lemma}
\label{lemma-descend-finite-type}
Let $X$ be a scheme. Let $\mathcal{F}$ be an $\mathcal{O}_X$-module.
The following are equivalent
\begin{enumerate}
\item $\mathcal{F}$ is of finite type as an $\mathcal{O}_X$-module, and
\item $\epsilon^*\mathcal{F}$ is of finite type as an
$\mathcal{O}_\etale$-module on the small \'etale site of $X$.
\end{enumerate}
Here $\epsilon$ is as in (\ref{equation-epsilon}).
\end{lemma}

\begin{proof}
The implication (1) $\Rightarrow$ (2) is a general fact, see
Modules on Sites, Lemma \ref{sites-modules-lemma-local-pullback}.
Assume (2). By assumption there exists an \'etale covering
$\{f_i : X_i \to X\}$ such that
$\epsilon^*\mathcal{F}|_{(X_i)_\etale}$ is generated by
finitely many sections. Let $x \in X$. We will show that $\mathcal{F}$
is generated by finitely many sections in a neighbourhood of $x$.
Say $x$ is in the image of $X_i \to X$ and denote $X' = X_i$. Let
$s_1, \ldots, s_n \in
\Gamma(X', \epsilon^*\mathcal{F}|_{X'_\etale})$
be generating sections. As
$\epsilon^*\mathcal{F} =
\epsilon^{-1}\mathcal{F} \otimes_{\epsilon^{-1}\mathcal{O}_X}
\mathcal{O}_\etale$
we can find an \'etale morphism $X'' \to X'$ such that $x$ is
in the image of $X' \to X$ and such that
$s_i|_{X''} = \sum s_{ij} \otimes a_{ij}$ for some sections
$s_{ij} \in \epsilon^{-1}\mathcal{F}(X'')$ and
$a_{ij} \in \mathcal{O}_\etale(X'')$. Denote $U \subset X$ the image
of $X'' \to X$. This is an open subscheme as $f'' : X'' \to X$ is \'etale
(Morphisms, Lemma \ref{morphisms-lemma-etale-open}). After possibly
shrinking $X''$ more we may assume $s_{ij}$ come from elements
$t_{ij} \in \mathcal{F}(U)$ as follows from the construction of
the inverse image functor $\epsilon^{-1}$. Now we claim that
$t_{ij}$ generate $\mathcal{F}|_U$ which finishes the proof
of the lemma. Namely, the corresponding map
$\mathcal{O}_U^{\oplus N} \to \mathcal{F}|_U$ has the property
that its pullback by $f''$ to $X''$ is surjective. Since $f'' : X'' \to U$
is a surjective flat morphism of schemes, this implies that
$\mathcal{O}_U^{\oplus N} \to \mathcal{F}|_U$ is surjective by
looking at stalks and using that
$\mathcal{O}_{U, f''(z)} \to \mathcal{O}_{X'', z}$
is faithfully flat for all $z \in X''$.
\end{proof}

\end{document}
